\section{Image generation}

\subsection{Purpose and workflow}

C-LARA-2 retains the image-generation workflow developed in the final versions
of C-LARA. Its purpose is not merely to attach an attractive image to each
page. A multi-page pedagogical text needs illustrations that fit the content
of individual pages while also forming a coherent visual world: recurring
people, objects and locations should remain recognisable, and linework,
colour, tone and degree of realism should not vary arbitrarily. The
availability of GPT-Image-1 to the project in 2025 made this workflow much
more practical than it had been with earlier image-generation models
\citep{RendinaEtAl2025}.

The workflow separates three reviewable stages. First, a human author supplies
a short style brief. The system expands it into a fuller description and
creates a sample image; the author can inspect, edit or regenerate both before
continuing. Second, the system proposes recurring visual \emph{elements}---for
example, people, objects and locations. The author can remove unsuitable
proposals or add missing ones, and the platform creates descriptions and
reference images for the selected elements. Finally, it creates one or more
candidate page images using the page text, approved style, relevant element
references and any page-level guidance. Independent element and page requests
can be run in parallel, but selection and acceptance remain human decisions.

\subsection{A two-page French example}

We illustrate the workflow with the small French dialogue already used in the
annotation discussion:

\begin{quote}
\textbf{Page 1:} «~Je t'aime~», dit-il.\\
\textbf{Page 2:} «~Vous vous moquez de moi~», dit-elle.
\end{quote}

The example is deliberately minimal. Its purpose is not to claim literary or
pedagogical depth, but to make the relationship among style, elements and
page-level generation visible. The two lines create a small unresolved story:
a declaration is followed by a response suggesting that the woman believes the
man is making fun of her. This gives the generator a local emotional task on
each page while requiring continuity of character and visual atmosphere.

For the style stage, the author supplied the brief ``Late nineteenth-century
French lithograph-inspired illustration, with restrained hand-coloured tones,
fine ink linework, soft paper texture, and a slightly melancholic literary
atmosphere.'' The brief also requested suitability for adult or older teenage
language learners and prohibited visible written text, speech balloons, labels
and signage. The platform expanded this in French into a more detailed prompt
specifying delicate ink lines, muted hand-coloured tones, paper texture, soft
light, subdued shadows and an intimate, contemplative composition. This is a
useful example of the division of labour: the human supplied the historical
and pedagogical intention, while the AI converted it into a fuller operational
specification that remained available for human review.

\subsection{Recurring elements and page images}

Element discovery proposed \emph{l'homme} and \emph{la femme}. Their expanded
descriptions specify nineteenth-century dress, fine ink detail, muted colour,
soft lighting and emotionally expressive but restrained faces. The resulting
reference images establish a recognisable visual identity for each person.
They are then provided as relevant element references for both page prompts,
so that the system has an explicit continuity mechanism rather than relying
only on repeated textual descriptions.

The page prompts are intentionally close to the element and style
specifications. Page 1 foregrounds the man and the emotionally tentative
moment of the declaration; Page 2 foregrounds the woman and her more guarded
response. Both prompts retain references to the woman and the man, even where
only one is foregrounded. Three variants were requested for each page, after
which a human author selected the preferred variant. This is an important
practical point: generation produces candidates, not an automatically accepted
final illustration.

% Repository figure files remain text-only placeholders. In Overleaf, replace
% them with a montage of the selected element-reference and page-image JPEGs.
% \begin{figure}[htbp]
% \centering
% \includegraphics[width=\linewidth]{images/lithograph_elements.jpg}
% \caption{Reference images for \emph{l'homme} and \emph{la femme}, generated
% from the approved late-nineteenth-century French lithograph style.}
% \label{fig:lithograph-elements}
% \end{figure}
%
% \begin{figure}[htbp]
% \centering
% \includegraphics[width=\linewidth]{images/lithograph_pages.jpg}
% \caption{Selected page-image variants for the two-page French dialogue. The
% common style and element references aim to preserve visual continuity while
% allowing the emotional focus to change between pages.}
% \label{fig:lithograph-pages}
% \end{figure}

\subsection{Human review and limitations}

The example illustrates why image generation is organised as a staged
workflow. A single page prompt can produce a plausible isolated picture, but
it gives little control over cross-page consistency. Separating style,
recurring elements and page images gives a human author several points at
which to inspect and redirect the process. It also makes failures easier to
locate: a problem may concern the style description, an element identity, a
page prompt or the selection among generated variants.

The workflow does not make visual or cultural judgement automatic. Generated
images can still misrepresent a relationship, omit a relevant detail, produce
inconsistent character appearances or be inappropriate for a learner or
community context. Visible text in images may also be actively undesirable,
particularly for picture-based exercises. For this reason, C-LARA-2 treats
image generation as AI-supported production under human review, with editing,
regeneration and selection available throughout. These concerns become still
more important for the specialised picture-dictionary workflow discussed in
the next section.
